\item A small disc \( A \) slides down with initial velocity equal to zero from the top of a smooth hill of height \( H \) having a horizontal portion (Fig. 1.30). What must be the height of the horizontal portion \( h \) to ensure the maximum distance \( s \) covered by the disc? What is it equal to?
    \begin{center}
        \begin{tikzpicture}
            \node at (0, 0) {Diagram};
        \end{tikzpicture}
    \end{center}
\begin{solution}
    \begin{center}
        \begin{tikzpicture}
            \pic at (0, 0) {frame=3cm};
        \end{tikzpicture}
    \end{center}
    
    \begin{align*}
        \intertext{Velocity of the body at height $h$, $v_b = \sqrt{2g(H - h)}$, horizontally (from the figure given in the problem book). Time taken in falling through the distance $h$ is given by}
        t &= \sqrt{\dfrac{2h}{g}} \quad \text{(as initial vertical component of the velocity is zero)}
        \intertext{Now,}
        s &= v_b t = \sqrt{2g(H + h)} \times \sqrt{\dfrac{2h}{g}} = \sqrt{4(Hh - h^2)}
        \intertext{For $s_{\text{max}}$,}
        \dfrac{d}{ds} \left(Hh - h^2\right) &= 0, \quad \text{which yields} \quad h = \dfrac{H}{2}
        \intertext{Putting this value of $h$ in the expression obtained for $s$, we get}
        s_{\text{max}} &= H
    \end{align*}
\end{solution}
